% !TEX root = 概率论期中复习.tex
\section{第八周：连续分布、对称性与无记忆性}

\subsection{教材第57页第2题}\label{ex:w08-b57-2}
设 $f,g$ 都是 $\R$ 上的概率分布密度函数，而 $p,q\in[0,1]$ 满足 $p+q=1$。证明 $pf+qg$ 是 $\R$ 上的概率分布密度函数。
\basis{原笔记第2章第3节命题“概率密度的基本性质”。}
\paragraph{思路} 分别核对非负性和归一化；这两项就是成为密度所需的条件。
\begin{proof}
记 $h=pf+qg$。由于 $p,q\geq0$ 且 $f,g$ 非负，有 $h(x)\geq0$。此外，$f,g$ 可积，故 $h$ 也可积，且
\begin{align*}
\int_{\R}h(x)\,dx
&=p\int_{\R}f(x)\,dx+q\int_{\R}g(x)\,dx\\
&=p+q=1.
\end{align*}
因此 $h$ 满足密度的非负性和归一化条件，是 $\R$ 上的概率分布密度函数。
\end{proof}
\paragraph{答案} $pf+qg$ 是密度；其分布函数为 $pF_f+qF_g$，其中 $F_f,F_g$ 分别是 $f,g$ 对应的分布函数。

\subsection{教材第58页第4题}\label{ex:w08-b58-4}
设 $\R$ 上的绝对连续型随机变量 $\xi$ 有密度函数
\[
f(x)=e^{-x-e^{-x}}.
\]
求 $\xi$ 的分布函数 $F$（称为极值分布）。
\basis{原笔记第2章第3节“密度与分布函数的关系”。}
\paragraph{思路} 将密度写成 $e^{-x}\exp(-e^{-x})$，用 $u=e^{-x}$ 换元。
\begin{proof}[解答]
对任意 $t\in\R$，
\begin{align*}
F(t)
&=\int_{-\infty}^t e^{-x}\exp(-e^{-x})\,dx\\
&=\int_{e^{-t}}^{\infty}e^{-u}\,du
=\exp(-e^{-t}).
\end{align*}
这里 $x\to-\infty$ 时 $u\to\infty$，且 $du=-e^{-x}\,dx$，因此换元后积分上下限及符号如上。所得函数在 $t\to-\infty$ 时趋于 $0$，在 $t\to\infty$ 时趋于 $1$。
\end{proof}
\paragraph{答案} $F(t)=\exp(-e^{-t})$，$t\in\R$。

\subsection{作业第2题：由三角形密度求分布函数}\label{ex:w08-h2}
设连续型随机变量 $X$ 的密度函数为
\[
f(x)=\begin{cases}x,&0\leq x\leq1,\\2-x,&1\leq x<2,\\0,&\text{其他。}\end{cases}
\]
求随机变量 $X$ 的分布函数 $F(t)$。
\basis{原笔记第2章第3节例“分段密度的参数、区间概率与分布函数”。}
\paragraph{思路} 按积分上限所在的密度分段分别计算；$t>1$ 时须包含第一段的全部积分。
\begin{proof}[解答]
$t<0$ 时，支撑集在 $t$ 的右侧，故 $F(t)=0$。$0\leq t\leq1$ 时，
\[
F(t)=\int_0^t x\,dx=\frac{t^2}2.
\]
$1<t<2$ 时，
\begin{align*}
F(t)&=\int_0^1x\,dx+\int_1^t(2-x)\,dx\\
&=\frac12+2t-\frac{t^2}2-\frac32
=1-\frac{(2-t)^2}2.
\end{align*}
$t\geq2$ 时已积尽全部概率，$F(t)=1$。两条中间表达在 $t=1$ 处均为 $1/2$，各分段相接连续。
\end{proof}
\paragraph{答案}
\[
F(t)=\begin{cases}0,&t<0,\\t^2/2,&0\leq t\leq1,\\1-(2-t)^2/2,&1<t<2,\\1,&t\geq2.\end{cases}
\]

\subsection{作业第3题：对称密度的分布函数}\label{ex:w08-h3}
设连续函数 $X$ 的密度函数 $f(x)$ 关于 $x=c$ 对称。求证：其分布函数 $F(t)$ 满足
\[
F(c-t)+F(c+t)=1,\qquad\forall t\in\R.
\]
注：不能只画图，要用关于密度的积分进行数学推导。
\basis{原笔记第2章第3节“密度与分布函数的关系”，命题“标准正态分布的对称公式”的积分思路。}
\paragraph{思路} 关于 $c$ 对称表示 $f(2c-x)=f(x)$，对左尾积分作反射换元即可得到右尾。
\begin{proof}
题面“连续函数 $X$”按具有所给密度的连续型随机变量理解。对称性给出 $f(2c-u)=f(u)$。对任意实数 $t$，在积分中令 $u=2c-x$，有
\begin{align*}
F(c-t)
&=\int_{-\infty}^{c-t}f(x)\,dx\\
&=\int_{c+t}^{\infty}f(2c-u)\,du\\
&=\int_{c+t}^{\infty}f(u)\,du
=1-F(c+t).
\end{align*}
最后一步使用密度的总积分为 $1$。将 $F(c+t)$ 移到等式左侧即得结论；这一换元对正、负及零的 $t$ 都成立。
\end{proof}
\paragraph{答案} $F(c-t)+F(c+t)=1$ 对全部 $t\in\R$ 成立；特别地 $F(c)=1/2$。

\subsection{作业第4题：先到列车的方向}\label{ex:w08-h4}
从早上7点开始，每隔15分钟都有一趟列车开往A地，而从7:05开始，每隔15分钟有一趟列车开往B地。
\begin{enumerate}
\item 如果某位乘客在7点到8点之间到达车站的时间服从均匀分布。他到达车站以后，无论下一趟进站的列车是开往A地还是开往B地的，他立刻上那一列车。那么他乘上开往A地的列车的概率是多大？
\item 如果该乘客在早上7:10到8:10之间到达车站的时间服从均匀分布，请回答(a)中相同的问题。
\end{enumerate}
\basis{原笔记第2章第3节“均匀分布”，例“等公交车不超过五分钟”。}
\paragraph{思路} 列出乘客会等到A车的到达时间段，以总长度除以均匀到达区间长度。
\begin{proof}[解答]
用从7点开始的分钟数表示时刻。A车在 $0,15,30,45,60,\ldots$ 到达，B车在 $5,20,35,50,65,\ldots$ 到达。乘客到达时刻恰为列车时刻的概率是 $0$，因此各段端点是否包含不影响答案。
\begin{enumerate}
\item 在 $[0,60]$ 内，乘客先等到A车的时间段为
\[
(5,15),\quad(20,30),\quad(35,45),\quad(50,60).
\]
总长度为 $4\times10=40$ 分钟。因此所求概率为 $40/60=2/3$。
\item 在 $[10,70]$ 内，相应时间段为
\[
(10,15),\quad(20,30),\quad(35,45),\quad(50,60),\quad(65,70).
\]
总长度为 $5+10+10+10+5=40$ 分钟，均匀到达区间长度仍为 $60$ 分钟，所以概率仍为 $2/3$。
\end{enumerate}
\end{proof}
\paragraph{答案} 两问均为 $2/3$。

\subsection{作业第5题：修理时间}\label{ex:w08-h5}
修理某机器所需的时间（单位：小时）服从参数为 $1/2$ 的指数分布。
\begin{enumerate}
\item 修理时间超过2小时的概率是多少？（结果保留3位小数）
\item 若已经持续修理了9小时的条件下，总共需要10小时才能修好的条件概率是多少？（结果保留3位小数）
\end{enumerate}
\basis{原笔记第2章第3节命题“指数分布的分布函数与尾概率”、定理“指数分布的无记忆性”。}
\paragraph{思路} 尾概率为 $e^{-t/2}$；第二问需区分“恰好10小时”和“10小时内修好”。
\begin{proof}[解答]
记总修理时间为 $T\sim\operatorname{Exp}(1/2)$。
\begin{enumerate}
\item 直接使用尾概率，
\[
\P(T>2)=e^{-2/2}=e^{-1}\approx0.368.
\]
\item 已持续修理9小时且尚未修好对应 $\{T>9\}$。若按题面字面“总共需要10小时”理解为恰好 $T=10$，连续型随机变量的单点概率为零，故
\[
\P(T=10\mid T>9)=\frac{\P(T=10)}{\P(T>9)}=0.000.
\]
若题意是“总修理时间不超过10小时”，则所求为
\begin{align*}
\P(T\leq10\mid T>9)
&=1-\P(T>10\mid T>9)\\
&=1-\frac{e^{-10/2}}{e^{-9/2}}
=1-e^{-1/2}\approx0.393.
\end{align*}
若题意是“至少还要修到第10小时”，相应的条件概率则为 $e^{-1/2}\approx0.607$。三者事件不同，不能混用。
\end{enumerate}
\end{proof}
\paragraph{答案} 第一问 $0.368$；第二问字面“恰好10小时”为 $0.000$，若指10小时内修好则为 $0.393$。

\subsection{作业第6题：旧车的剩余里程}\label{ex:w08-h6}
甲估计，一辆汽车在报废之前能行驶的里程数（单位：千里）服从参数为 $1/20$ 的指数分布。乙有一辆已经行驶了10公里的旧车。如果甲买下此车，那么该车至少还能行驶20千里的概率是多少？（结果保留3位小数）如果将以上的指数分布换为 $(0,40)$ 上的均匀分布，再求以上概率。
\basis{原笔记第2章第3节“指数分布的无记忆性”“均匀分布”。}
\paragraph{思路} 先统一单位，以“总里程已超过旧车里程”为条件计算剩余寿命的尾概率。
\begin{proof}[解答]
令 $L$ 表示总里程，以千里为单位，令 $d$ 表示旧车已经行驶的里程换算为相同单位后的数值。剩余里程至少为 $20$ 的条件概率是
\[
\P(L\geq d+20\mid L>d).
\]
对指数分布，只需 $d\geq0$，由无记忆性有
\[
\P(L\geq d+20\mid L>d)
=\frac{e^{-(d+20)/20}}{e^{-d/20}}
=e^{-1}\approx0.368.
\]
对 $L\sim U(0,40)$，当 $0\leq d<20$ 时，
\[
\P(L\geq d+20\mid L>d)
=\frac{40-d-20}{40-d}=\frac{20-d}{40-d}.
\]
题面总里程单位“千里”与已行驶“10公里”并不相同。按“里”为市里，$10$ 公里等于 $20$ 里，即 $d=0.02$ 千里，此时
\[
\frac{20-0.02}{40-0.02}=\frac{19.98}{39.98}\approx0.500.
\]
若题面“千里”实际用来指千公里，则 $d=0.01$，概率为 $19.99/39.99\approx0.500$。若原意是“已经行驶10千里”，则 $d=10$，均匀分布下的答案改为 $10/30\approx0.333$。指数分布的答案在这些解释下都为 $0.368$。
\end{proof}
\paragraph{答案} 指数分布下为 $0.368$；保留题面“10公里”并统一单位后，均匀分布下为 $0.500$；若原意为“10千里”，则为 $0.333$。

\subsection{作业第7题：正态分布的中心}\label{ex:w08-h7}
设 $X\sim N(\mu,4)$，并且 $\P(X\geq3)=0.5$。求 $\mu$ 是多少？
\basis{原笔记第2章第3节“标准正态分布与对称性”“正态分布的线性变换与标准化”。}
\paragraph{思路} 正态分布在其均值两侧各有一半概率，并利用其分布函数严格递增确定唯一的中心。
\begin{proof}[解答]
标准差为 $2$。令 $Z=(X-\mu)/2\sim N(0,1)$，以 $\Phi$ 表示标准正态分布函数，则
\[
\P(X\geq3)=1-\Phi\left(\frac{3-\mu}{2}\right)=\frac12.
\]
所以 $\Phi((3-\mu)/2)=1/2$。由于 $\Phi(0)=1/2$，且标准正态密度处处正，$\Phi$ 严格递增，得到 $(3-\mu)/2=0$，即 $\mu=3$。
\end{proof}
\paragraph{答案} $\mu=3$。
